\section{Results}
\label{sec:results}

The results divide along the paper's claim. Sections~\ref{sec:compute1} to
\ref{sec:compute2} concern computing over the graph's structure, where it pays. The final
part concerns showing the graph to a model, where it does not.

\subsection{Detection: what declared signals catch}

A first corpus crashed on seventeen of twenty-eight faults, so execution is itself a strong
free detector; but a fault that halts the run never reaches a trust decision, and real
generated pipelines read inputs through accessors with defaults. Everything below concerns
faults such a pipeline survives silently: 104 cases across two pipelines, the cross product
of five fault classes, the boundaries and the applicable parameters, most from classes
observed in this system's own logs.

Structural schema validation leaves 78\,\% undetected, declared invariants 27\,\%, and a
model judge 18\,\%, none of them raising a single false suspicion. The intervals for the
latter two overlap and an exact McNemar test does not separate them ($p=0.16$), so we make
no claim that either is better; what the sample supports is that both are significantly
worse than their union, which leaves 7\,\% ($p<0.0001$ and $p=0.0005$). The failures are
disjoint. Invariants catch twelve faults the judge misses, every one a matter of counting
or shape; the judge catches twenty-one the invariants miss, eighteen of them substituted
identifiers including every plausible one that defeated the pattern check. Schema
validation contributes nothing either does not, detecting only the fault class that
changes a record's shape.

\subsection{What the graph is actually for}
\label{sec:compute1}

Detection tells a run something is wrong. It does not say what to repair. When several
boundaries are flagged, something must choose, and that choice is where the graph does work
no cheaper signal can. We measure a selector that walks recorded dataflow to the causal
root against one that walks the declared boundary list in order.

The comparison is only meaningful when those two orders differ. In a pipeline whose
functions are declared in the order they run, the causal root of any flagged set is simply
the first of them, and the selectors cannot disagree. Models writing code do not reliably
declare in execution order, so we test both conditions across four shapes: a deep chain, a
wide fan-out, a diamond that splits and reconverges, and a short chain.

\begin{table}[t]
\centering
\begin{tabular}{llrrr}
\toprule
Shape & Declaration order & Contested & Causal root & Source order \\
\midrule
chain-8    & dataflow  & 16 & 100\,\% & 100\,\% \\
chain-8    & scrambled & 16 & 100\,\% & 38\,\% \\
fan-6      & scrambled &  2 & 100\,\% & 0\,\% \\
diamond-4  & scrambled & 14 & 100\,\% & 14\,\% \\
chain-3    & scrambled &  6 & 100\,\% & 0\,\% \\
\midrule
all        & dataflow  & 38 & 100\,\% & 100\,\% \\
all        & scrambled & 38 & \textbf{100\,\%} & \textbf{21\,\%} \\
\bottomrule
\end{tabular}
\caption{Repair-target selection across four pipeline shapes. The causal-root selector is
exact in every cell. The baseline matches it only when declaration order already encodes
execution order.}
\label{tab:selection}
\end{table}

The causal-root selector is exact in all 38 contested cases of every shape and both
orderings. The baseline matches it perfectly when declaration order happens to encode
execution order and falls to 21\,\% when it does not. This gives the graph's value a precise
form: \emph{it is worth exactly the extent to which declaration order diverges from
execution order}. A team whose pipelines are written in order gains nothing here. A team
whose pipelines are model-authored gains the difference between 100 and 21 per cent, and
cannot know in advance which case they are in, which is itself an argument for computing
the relation rather than assuming it.

\subsection{The model judge, and where each signal fails}
\label{sec:compute2}

We expected the language-model judge to be the weak arm. It was not, and reporting what
happened instead is more useful than the result we set out to find.

The judge is a 3B model run locally, asked about one boundary at a time and given that
stage's declared purpose in prose. Across both corpora, 104 silent faults with three
repeats each, it never once disagreed with itself. Every one of the 104 cases came back
unanimous over 312 calls, with no parse failures.

A false-trust rate alone would mean nothing, because a judge answering suspect
indiscriminately scores perfectly while being useless. We therefore asked the same judge
about boundaries known to be sound: every boundary of the clean unfaulted pipeline, and
for each case a boundary strictly upstream of the fault, which cannot be contaminated.
Across 102 such calls it returned trusted every time, so its false-suspicion rate matches
both deterministic signals at zero and the comparison stands.

\begin{table}[t]
\centering
\begin{tabular}{lrrl}
\toprule
Signal & Missed & Rate & 95\,\% CI \\
\midrule
schema validation      & 81/104 & 78\,\% & [69, 85] \\
declared invariants    & 28/104 & 27\,\% & [19, 36] \\
model judge (3B)       & 19/104 & 18\,\% & [12, 27] \\
invariants + judge     &  7/104 &  7\,\% & [3, 13] \\
\bottomrule
\end{tabular}
\caption{Faults left undetected over 104 cases, with Wilson intervals. The intervals for
invariants and the judge overlap; the pair is separated from both.}
\label{tab:signals}
\end{table}

\subsection{What the sample supports, and what it does not}

The ranking in Table~\ref{tab:signals} invites a conclusion the data does not carry, so
we tested it on the paired outcomes. An exact McNemar test does not separate the judge
from declared invariants (12 against 21 discordant, $p=0.16$), and their confidence
intervals overlap. We make no claim that either is the better signal, and the sample is
large enough that this is a finding rather than a shrug: two signals of quite different
kinds, one a compiled check and one a 3B model, are not distinguishable by error rate.

The combination is a different matter. It improves on invariants alone (0 against 21,
$p<0.0001$) and on the judge alone (0 against 12, $p=0.0005$). Both single signals are
significantly worse than using both; neither is significantly worse than the other.
Schema validation is separated from all three at $p<0.0001$, and adding it to any pair
changes nothing, because it catches nothing they do not.

The result is therefore about complementarity rather than ranking. Invariants catch
twelve faults the judge misses, and every one is a matter of counting or shape: fields
removed from a record, sequences silently shortened. The judge catches twenty-one the
invariants miss, and eighteen of those are substituted identifiers, including every
plausible one that defeated the pattern check. Each signal is blind close to where the
other sees, and it is that disjointness the evidence establishes.

It would be wrong to read this as a model outperforming a static check. The judge was
given the stage's declared purpose, for instance that a retrieval stage must fetch every
source in every declared class and report what it retrieved. That is an acceptance
criterion. It is the same kind of statement the invariant encodes, written in prose and
checked by a model rather than compiled into a count and a regular expression.

So the comparison was never between having criteria and not having them. It was between
two encodings of the same criteria, and the encodings fail differently. A count is exact
about quantity and silent about meaning. A sentence carries meaning and is vague about
quantity. That is why the union is so much stronger than either part, and it is a reason
to declare criteria at a boundary at all rather than a reason to prefer one checker.

\subsection{Showing the graph to a model does not pay}
\label{sec:show}

The use an execution graph is most often justified by is presentation: give the agent a
relevant slice rather than everything or nothing. Four regimes were asked the same
question over a branching corpus, name the faulty stage, with the injected fault as the
answer.

\begin{table}[t]
\centering
\begin{tabular}{lrrl}
\toprule
Regime & Accuracy & Context & Note \\
\midrule
stage purposes only        & 15\,\% & 1{,}343 & no captured values \\
size-matched random slice  & 35\,\% & 1{,}402 & control \\
graph-selected region      & 40\,\% & 1{,}466 & $p=0.69$ vs random \\
selected plus expansion    & 46\,\% & 1{,}466 & the deployed method \\
complete current run       & 50\,\% & 1{,}913 & \\
accumulated history, d=64  & 40\,\% & 129{,}246 & \\
\bottomrule
\end{tabular}
\caption{Localisation under six context regimes. Selection is not separable from a random
slice of the same size, including in the form the system actually deploys; accumulated
history costs two orders of magnitude more for less.}
\label{tab:context}
\end{table}

Captured values matter enormously, 15\,\% against 40 to 50 for any regime showing what the
code produced, which is the strongest evidence here that execution records are the right
substrate. Selecting among them does not help: the graph-selected region is not separable
from a size-matched random slice ($p=0.69$), and adding expansion on request, which is the
method as deployed rather than a weakened stand-in, reaches 46\,\% against 50\,\% for
simply showing everything. Accumulated history degrades from 52\,\% at depth four or less
to 40\,\% at depth sixteen or more while its context grows to 129{,}246 characters against
1{,}913. The mechanism is visible in the failures: shown captured values, a model names a
downstream stage in 83\,\% of cases, and selection acts on the prompt rather than on the
answer, so a model handed a selected region still blames the most downstream thing in it.

\subsection{A protocol that withholds and is asked}
\label{sec:collapsed}

The arms above hand the reviewer a slice chosen in advance. The system we study does
something different, and the distinction turns out to matter. Its initial package carries
nodes at exactly the boundary's function-stack level, names the nested helpers without
their contents, and permits the reviewer to request one direct child at a time for at most
three rounds. Evaluating that needs a pipeline whose boundaries call helpers, so we built
one and planted faults at three depths: visible at the boundary, inside a helper, and
inside a sub-helper.

\begin{table}[t]
\centering
\begin{tabular}{lrrrr}
\toprule
Arm & Faults & Clean & Mean chars & Expansions \\
\midrule
collapsed, expansion on request & 5/6 & 4/4 & 841 & 9 \\
everything, flattened           & 5/6 & 4/4 & 922 & 0 \\
random subset, budget matched   & 4/6 & 3/4 & 831 & 0 \\
\bottomrule
\end{tabular}
\caption{The deployed protocol against sending everything and against a size-matched
random subset. It matches the former on nine per cent less context and, unlike every
fixed-slice arm we tested, does not tie the latter.}
\label{tab:collapsed}
\end{table}

This is the only selection arm in the paper that separates from a size-matched random
control, and it forces a narrower statement of the result above. A selected region handed
over in advance is not better than a random one of the same size. A protocol that
withholds nested detail and lets the reviewer pull what it needs is a different mechanism,
and on this evidence it works: the same accuracy as sending everything, for less, and
better than random at equal cost. What we can say is that the benefit lies in the
interaction rather than in the choice of slice, which is not what we expected and not what
the fixed-slice arms could have shown.

Two details qualify it. Faults planted two levels deep were caught with zero and one
expansions, because a deep fault usually surfaces in the boundary's own output, so the
reviewer rarely has to drill to where the fault lives. And the single miss is a reversed
ranking judged clean after an expansion was spent, which is precisely the fault an
ordering scar catches deterministically in Section~\ref{sec:scars}. With ten cases per arm
this is directional rather than sized, and we implemented the collapsed-helper half of the
protocol but not its artifact inspection.

\subsection{External validity: three recovered runs of the real system}
\label{sec:external}

Everything above runs on pipelines we authored, which is the weakest thing about it. Three
records of the real system, deleted in a cleanup commit but recoverable from its git
history, supply the external check.

A four-arm run of 29 July 2026, on the real model-authored pipeline with a frontier model,
reports that graph-selected context matched the full graph's issue coverage with 46\,\%
fewer input tokens on the initial run and 47\,\% on the latest, while accumulated flat
history grew from 71{,}470 to 303{,}757 characters and became 29\,\% more expensive than
the selected region. A second run of 30 July shows the same shape at 31{,}369 tokens for
stage metadata against 109{,}590 for accumulated history, all arms finding the same three
issue functions. This is the efficiency half of our claim, measured on the system itself.

The same 29 July record contains one further observation we regard as the most important
sentence in either dataset: accumulated history \emph{retained a previously repaired
boundary as an open issue after every other arm treated it as resolved}, which its authors
attribute to stale prior versions remaining salient in flat history. Our synthetic history
sweep never produced this, and the recovered run shows why: our accumulated history was
clean prior iterations, which hand the model a helpful baseline to diff against, whereas
real history contains superseded versions of boundaries that have since changed. Staleness,
not volume, is the active ingredient. Section~\ref{sec:stale} tests that as an
intervention.

\subsection{Staleness, briefly}
\label{sec:stale}

Accumulated history harms a model judge, and what it does depends on the model. A
superseded faulty run in context halves one judge's detection of genuine new faults;
another judge cannot recognise a clean run without one, inventing a fault in sixteen of
twenty clean cases and naming the final stage every time; the deployed system's own record
shows a third model resurrecting a resolved issue from flat history. Neither of ours
re-flagged the repaired boundary, which was the predicted failure. The stable conclusion
is a dependence rather than a direction: history content modulates judge behaviour in
model-specific ways that no prompt here controlled, while deterministic signals are
unaffected by context by construction. Full cells in Appendix~\ref{app:stale}; the
detection drop does not survive the multiplicity correction of Section~\ref{sec:power}.

\subsection{Checks placed where faults enter}
\label{sec:scars}

Every signal so far was written in advance by an author deciding what to assert. The
alternative is to learn a check from a fault once it has been caught, place it on the
boundary where it entered, and pay for that fault only once.

We adopt the term \emph{scar} from prior work on coupling kernels, where a scar is a
low-rank invariant subspace that lets a query be \emph{read} exactly rather than searched
over an intractably large state space \cite{rovai2026worldkernel}, and from its use in
agent goal induction, where a goal is treated as a low-rank commitment set read in a
number of probes proportional to its rank \cite{rovai2026civex}. The move transfers
directly. A caught fault commits a small set of properties; rather than re-deriving them
on every subsequent run, they are read off once and pinned to the boundary that produced
them. What is new here is the application to trust decisions and the measurement of what
it buys, which separates the two halves of this paper more sharply than anything else.

A scar is minted from a population of healthy observations of one boundary and one faulty
observation of the same boundary, keeping only structural properties that held across the
whole population and broke in the faulty one: a field present throughout and now absent, a
count below its healthy floor, a list that stopped being homogeneous, an ordering that
held and no longer does. A property expressible as no such check is refused rather than
stored as vague suspicion, because the deterministic signals' zero false-suspicion rate is
the thing that makes them worth having.

\begin{table}[t]
\centering
\begin{tabular}{lrrr}
\toprule
Signal & Detected & Localised & Misplaced \\
\midrule
declared contracts & 26/90 & \textbf{0/90} & 26/90 \\
learned scars      & 60/90 & \textbf{60/90} & 0/90 \\
\bottomrule
\end{tabular}
\caption{Ninety data faults. Localisation is whether the flagged boundary is the one the
fault entered. Declared contracts never localise: every detection fires on an aggregating
stage downstream of the cause.}
\label{tab:scars}
\end{table}

Table~\ref{tab:scars} contains the result we consider most important in this paper.
Hand-written contracts detect 26 of 90 faults and localise \emph{none} of them. A stage
that faithfully passes on a bad input satisfies every criterion anyone thought to write
about its own behaviour, so the checks fire on the aggregating stages downstream and the
suspect set sits entirely past the cause. This is the same failure the model judges
showed, in a signal with no model in it, which tells us the downstream-blame bias is not a
property of language models but of where checks are placed. Scars localise every fault
they detect, because a scar is placed where the fault was observed to enter rather than
where an author's attention happened to fall.

Two further results bound the method. Population size does not affect detection at all,
identical at one, four and eight healthy observations; what it buys is false suspicion,
falling from 21 of 56 boundary-checks at $n=1$ to zero at $n=8$. A single healthy
observation cannot separate an invariant from that run's value, so scars minted from one
pair fire on honest variation, and only a population can tell the difference. Second, a
vocabulary check is minted only when the population shows the domain closed: identifiers
drawn from a fixed registry stop admitting new values, dates and free text do not, and a
vocabulary check over an open domain is guaranteed to fire on an honest run eventually.
Enforcing that costs detection, 73 falling to 60, and takes false suspicion on held-out
healthy runs the scars never saw from 3 of 28 to zero. We regard that trade as the correct
one, and the refusal of faults visible only in an open domain as a feature of the design.

We also implemented the two graph computations this suggests, and both are negative.
Forward propagation of contamination reclassifies nothing, because the declared checks
fire downstream of the fault and there is nothing to propagate forward from. Backward
attribution narrows to seven boundaries of seven, because in a pipeline that funnels
through a merge every branch is an ancestor of the first suspect. Both failures have one
cause: exoneration requires a check that the fault \emph{would have violated}, not merely
a check that passed, and a hand-written contract does not supply one where a scar does.
